\documentclass[10pt,a4paper]{article}

\usepackage[a4paper,top=0.65cm,bottom=0.65cm,left=0.9cm,right=0.9cm]{geometry}
\usepackage[utf8]{inputenc}
\usepackage[T1]{fontenc}
\usepackage{lmodern}
\input{glyphtounicode}
\pdfgentounicode=1
\usepackage[french]{babel}
\usepackage[hidelinks]{hyperref}
\usepackage{xcolor}
\usepackage{titlesec}
\usepackage{enumitem}

\pagestyle{empty}
\setlength{\parindent}{0pt}
\setlength{\parskip}{0pt}
\linespread{0.92}
\setlist[itemize]{leftmargin=1.05em,itemsep=0pt,topsep=1pt,parsep=0pt,label=-}
\urlstyle{same}

\definecolor{heading}{HTML}{123867}
\titleformat{\section}{\normalsize\bfseries\color{heading}}{}{0pt}{}[\vspace{-5pt}\rule{\linewidth}{0.5pt}]
\titlespacing{\section}{0pt}{4pt}{2pt}

\newcommand{\role}[3]{%
  \textbf{#1} \hfill #3\\
  \textit{#2}\vspace{1pt}
}

\newcommand{\edu}[4]{%
  \textbf{#1} -- #2 \hfill #3 #4\\[-1pt]
}

\begin{document}

{\Large \textbf{Kossi Richard Allado}}\\
\textbf{Étudiant ingénieur IA et cybersécurité | IA générative, RAG, génie logiciel}\\
Beni Mellal, Maroc | +212 621 074 305 | \href{mailto:alladokossirichard2025@gmail.com}{alladokossirichard2025@gmail.com}\\
\href{https://allakori.github.io/}{allakori.github.io} | \href{https://github.com/ALLAKORI}{github.com/ALLAKORI} | \href{https://www.linkedin.com/in/kossi-richard-allado-50177b2a5}{linkedin.com/in/kossi-richard-allado}

\section{Profil}
Étudiant ingénieur en cybersécurité et intelligence artificielle à l'ENSA Beni Mellal, intéressé par un PFA en \textbf{IA appliquée à l'ingénierie logicielle}. Expérience pratique en Python, IA générative, RAG, analyse de données, modélisation fonctionnelle, documentation technique et prototypage full-stack. Motivé par les cas d'usage concrets : structuration de besoins, assistance IA, amélioration de prompts, classification d'informations, intégration de modèles et qualité logicielle.

\section{Compétences clés pour Univade}
\textbf{IA générative / RAG :} LLMs, prompt engineering, embeddings, Qdrant, FastEmbed, GroqCloud, LLaMA, Chainlit, LangChain basics, streaming SSE.\\
\textbf{Analyse et modélisation :} analyse de besoins, structuration d'informations métier, UML, modélisation de données, documentation fonctionnelle et technique.\\
\textbf{Data / Classification :} Python, Pandas, Scikit-learn, préparation de données, classification supervisée, Random Forest, évaluation de modèles.\\
\textbf{Génie logiciel :} FastAPI, Next.js, TypeScript, SQL, PostgreSQL, Redis, Docker, Git/GitHub, tests backend, API REST.\\
\textbf{Qualité et sécurité :} JWT, RBAC, rate limiting, masquage PII, audit logs, revue structurée, sens du détail et rigueur.

\section{Projets pertinents}
\textbf{USMS Chatbot - Assistant académique RAG} \hfill IA générative / RAG / Full-stack\\
\href{https://github.com/ALLAKORI/Chatbot-USMS}{github.com/ALLAKORI/Chatbot-USMS}
\begin{itemize}
  \item Conçu un assistant intelligent permettant l'accès en langage naturel à des FAQ, emplois du temps et informations administratives.
  \item Mis en place une architecture RAG avec FastAPI, Next.js, PostgreSQL, Redis, Qdrant, FastEmbed embeddings et génération de réponses via GroqCloud / LLaMA.
  \item Ajouté support multilingue, adaptation du ton, streaming SSE, tableau d'administration, authentification JWT, RBAC, limitation de requêtes, masquage PII, audit logs et tests backend.
\end{itemize}

\textbf{Expérimentations LLM local et chatbot} \hfill Prompting / Agents conversationnels\\
\href{https://github.com/ALLAKORI/MY-FIRST-AI-CHATBOT}{github.com/ALLAKORI/MY-FIRST-AI-CHATBOT}
\begin{itemize}
  \item Prototypé un chatbot local avec Python, Chainlit, LangChain basics, CTransformers et modèles GGUF.
  \item Testé la formulation de prompts, la logique conversationnelle et les limites des modèles locaux pour améliorer la qualité des réponses.
\end{itemize}

\textbf{High-Pressure Pump Fault Prediction} \hfill Machine Learning / Classification\\
\href{https://github.com/ALLAKORI/Defaut-Pompe-Prediction}{github.com/ALLAKORI/Defaut-Pompe-Prediction}
\begin{itemize}
  \item Développé un prototype de maintenance prédictive à partir de mesures capteurs : tension, puissance, vitesse, températures, pression et débit.
  \item Entraîné un modèle Random Forest avec Python, Pandas et Scikit-learn, puis créé une interface Tkinter pour tester le niveau de risque.
\end{itemize}

\textbf{Portfolio technique IA et cybersécurité} \hfill Documentation / Travail structuré\\
\href{https://allakori.github.io/}{allakori.github.io}
\begin{itemize}
  \item Organisé des projets sous forme de cas d'étude : objectif, architecture, outils, méthode, résultats, limites et pistes d'amélioration.
  \item Rédigé des write-ups techniques avec preuves, requêtes, chronologies et explications orientées qualité.
\end{itemize}

\section{Expérience}
\role{Stagiaire IA - Maintenance prédictive}{ONEE - Branche Eau, Kasba Tadla}{Août 2025}
\begin{itemize}
  \item Préparé et analysé des données capteurs de pompe haute pression pour prédire des défauts potentiels.
  \item Construit un prototype Python combinant modèle Random Forest, pipeline de préparation de données et interface de test.
\end{itemize}

\section{Formation}
\edu{Cycle ingénieur}{Cybersécurité et Intelligence Artificielle, ENSA Beni Mellal}{2024 - Présent}{}
\edu{DEUST}{Mathématiques, Informatique, Physique, FST Fès}{2022 - 2024}{-- Mention Bien}
\edu{Baccalauréat scientifique}{Série D, Lycée de Djagblé, Togo}{2022}{-- Mention Très Bien}

\section{Certifications et réalisations}
IBM Introduction to Data Engineering | Cisco Ethical Hacker | Google Cybersecurity: Foundations and Play It Safe | TryHackMe Junior Penetration Tester Path | CSIA CTF 2026 : 1ère place avec Team GHOSTSHELL

\section{Langues}
Français : courant | Anglais : professionnel technique | Éwé : langue maternelle

\end{document}
